\chapter{JSON API, CORS és böngészős kliensintegráció}

Az RWLang nem kényszerít választásra a klasszikus szerver-renderelt HTML és a JSON API között. Ugyanabban az alkalmazásban lehetnek \code{Html}-t visszaadó oldalak, formot feldolgozó actionök és typed JSON endpointok. Ez hasznos akkor is, ha csak egy kisebb JavaScript komponens kér adatot, és akkor is, ha külön frontend alkalmazás használja az RWLang szervert API-ként.

A fejezet célja három működési modell tiszta szétválasztása:

\begin{enumerate}
  \item \textbf{same-origin HTML + API}: a böngésző ugyanarról az originről tölti az oldalt és hívja az API-t;
  \item \textbf{cross-origin publikus API}: külön origin hívja az API-t, de nincs session-cookie;
  \item \textbf{cross-origin session auth}: külön frontend origin használ cookie-alapú sessiont, ezért CORS és CSRF policy együtt szükséges.
\end{enumerate}

A legfontosabb alapelv: \textbf{a CORS nem autentikáció és nem CSRF-védelem}. A CORS böngészőoldali hozzáférési policy; az auth azt dönti el, ki a felhasználó, a CSRF pedig azt védi, hogy egy autentikált session ne legyen más originről jogosulatlanul felhasználható állapotmódosításra.

\section{Az első typed JSON endpoint}

Egy JSON endpoint handlerének visszatérési típusa \code{Json}. A \code{json(...)} a támogatott runtime értéket JSON reprezentációvá alakítja.

\begin{lstlisting}[language=RWLang]
page fn statusApi(ctx: PageContext) -> Result<Json, PageError> {
    let healthy = true;
    return Ok(json(healthy));
}

route statusApi GET "/api/status" => statusApi;
\end{lstlisting}

A sikeres válasz reprezentációja JSON, például:

\Needspace{9\baselineskip}
\begin{lstlisting}
HTTP/1.1 200 OK
Content-Type: application/json; charset=utf-8

true
\end{lstlisting}

A \code{json(value)} nem tetszőleges objektum-serializáló. A támogatott értékek kötöttek: scalar értékek, model record, optional model és modell-lista. Ez illeszkedik az RWLang typed boundary szemléletéhez: a response alakja nem dinamikus, requestenként változó adatszerkezetből épül fel.

\section{Adatbázisból JSON lista}

A typed SQL és a JSON response közvetlenül összekapcsolható. A query továbbra is pontos model shape-et ad vissza, a JSON réteg pedig ezt reprezentálja.

\begin{lstlisting}[language=RWLang]
model Product {
    id: Int
    name: String
    price: Int
}

query fn listProducts(db: Db, limit: Int, offset: Int)
    -> Result<List<Product>, DbError> sql {
    SELECT id, name, price
    FROM products
    ORDER BY id
    LIMIT :limit OFFSET :offset
}

page fn productsApi(ctx: PageContext, db: Db,
                    limit: Int, offset: Int)
    -> Result<Json, PageError> {
    let products = listProducts(db, limit, offset)?;
    return Ok(json(products));
}

route productsApi GET "/api/products"
    query limit<Int> offset<Int>
    validate limit range 1 100 offset range 0 1000000
    => productsApi;
\end{lstlisting}

A fontos rész itt nem a JSON, hanem az, hogy az input ugyanúgy typed és validált, mint HTML route esetén. A kliens nem tudja például egy korlátlan \code{limit} értékkel megkerülni az alkalmazás route policy-ját.

\Needspace{17\baselineskip}
\section{Typed JSON POST body}

JSON request bodyhoz a route \code{json} input módot deklarál.

\begin{lstlisting}[language=RWLang]
action fn echoApi(ctx: ActionContext,
                  name: String,
                  age: Int,
                  active: Bool)
    -> Result<Json, PageError> {
    return Ok(json(name));
}

route echoApi POST "/api/echo"
    json name<String> age<Int> active<Bool>
    validate name length 1 100 age range 0 150
    auth user
    => echoApi;
\end{lstlisting}

Egy megfelelő kérés:

\begin{lstlisting}
POST /api/echo HTTP/1.1
Content-Type: application/json
Accept: application/json
X-CSRF-Token: <session CSRF token>

{"name":"Alice","age":42,"active":true}
\end{lstlisting}

A JSON input schema \textbf{zárt}. A runtime hibának tekinti többek között az ismeretlen vagy hiányzó kulcsot, a duplikált kulcsot, a rossz scalar típust, valamint a nem támogatott nested object, array, float vagy \code{null} inputot. Ezek request-contract hibák, ezért \code{400 Bad Request} választ adnak. Hibás vagy hiányzó \code{Content-Type} JSON route-on \code{415 Unsupported Media Type}. A mező- és request-méretlimitek ugyanúgy érvényesek, mint más request body esetén.

A \code{json}, \code{form} és \code{upload} body módok ugyanazon POST route-on kölcsönösen kizárják egymást, és ezt már a compiler ellenőrzi. JSON route \code{application/json}, egyszerű form route \code{application/x-www-form-urlencoded}, upload route pedig \code{multipart/form-data} requestet vár. Egy endpointnak legyen egyértelmű request contractja.

\section{Content negotiation és a 406 válasz}

A kliens az \code{Accept} headerrel jelezheti, milyen reprezentációt tud fogadni. JSON route-nál támogatott például:

\begin{lstlisting}
Accept: application/json
Accept: application/*
Accept: */*
\end{lstlisting}

Ha a handler JSON választ adna, de a kliens kizárólag \code{text/html} típust fogad el, a szerver \code{406 Not Acceptable} választ ad. Ez jobb annál, mintha a szerver csendben olyan reprezentációt küldene, amelyet a kliens explicit elutasított.

\section{API hibák}

JSON route-on a runtime által klasszifikált alkalmazáshibák ugyanazt az egyszerű envelope-ot használják: egyetlen stabil \code{error} kódot adnak vissza, miközben a részletes diagnosztika a szerveroldali logban marad. A kliens ne angol hibaüzenetek szövegét parse-olja, hanem a HTTP státuszt és ezt a kódot kezelje.

JSON-return route felismerése után a hibák JSON error envelope-ban érkeznek. Például:

\begin{lstlisting}
HTTP/1.1 401 Unauthorized
Content-Type: application/json; charset=utf-8

{"error":"unauthorized"}
\end{lstlisting}

Tipikus error kódok:

\begin{itemize}
  \item \code{bad_request};
  \item \code{unauthorized};
  \item \code{forbidden};
  \item \code{csrf_failed};
  \item \code{unsupported_media_type};
  \item \code{database_unavailable};
  \item \code{resource_limit};
  \item \code{internal_error}.
\end{itemize}

Ez különösen fontos böngészős kliensnél. Egy autentikált JSON route névtelen kérésre \code{401}-et ad, nem HTML login redirectet. A frontend így státuszkód alapján tud dönteni, nem kell HTML választ API hibaként értelmeznie.

\section{Same-origin: az alapértelmezett és legegyszerűbb modell}

Ha a HTML és az API ugyanazon scheme + host + port kombinációról érhető el, a böngésző szempontjából same-origin alkalmazásról beszélünk.

\begin{lstlisting}
https://app.example/
https://app.example/api/products
\end{lstlisting}

Ilyenkor általában nincs szükség CORS konfigurációra. Egy egyszerű kliens:

\begin{lstlisting}
const response = await fetch("/api/products?limit=20&offset=0", {
  headers: { Accept: "application/json" }
});

if (!response.ok) {
  throw new Error(`HTTP ${response.status}`);
}

const products = await response.json();
\end{lstlisting}

Ez legyen az alapértelmezett architektúra akkor, ha nincs valós üzleti igény külön frontend originre. Kevesebb trust boundaryt, egyszerűbb cookie policy-t és kevesebb deployment-konfigurációt jelent.

\section{CORS: exact origin allowlist}

Cross-origin böngészős hozzáférés alapból nincs engedélyezve. A trusted server configban konkrét origineket kell felsorolni; wildcard origin nincs.

\begin{lstlisting}
[web]
cors_origins = ["https://frontend.example"]
cors_allow_credentials = false
\end{lstlisting}

A CLI megfelelője:

\begin{lstlisting}
rwlang-server --cors-origin https://frontend.example ...
\end{lstlisting}

Több originhez több explicit allowlist elem tartozzon. Ne építs olyan szabályt, amely requestből származó \code{Origin} értéket vakon visszatükröz.

\section{Preflight}

Bizonyos cross-origin requestek előtt a böngésző \code{OPTIONS} preflight kérést küld. Az RWLang csak akkor enged preflightot, ha:

\begin{itemize}
  \item az \code{Origin} konfigurált;
  \item a cél GET/POST route létezik;
  \item a kért request header támogatott.
\end{itemize}

A támogatott browser API headerek között szerepel a \code{Content-Type}, az \code{Accept} és az \code{X-CSRF-Token}. A preflight tehát nem általános ``mindent engedő'' OPTIONS endpoint, hanem a tényleges route- és origin-policy része.

\section{Credentialed CORS session-cookie-val}

Ha a frontend külön originről használ RWLang session authot, credentialed CORS szükséges:

\begin{lstlisting}
[web]
cors_origins = ["https://frontend.example"]
cors_allow_credentials = true
\end{lstlisting}

Productionben ez csak HTTPS mellett engedélyezett. A session cookie ilyen konfigurációban \code{SameSite=None; Secure} tulajdonságot kap, mert a böngészőnek cross-site requestben is el kell küldenie.

A JavaScript oldalon a fetch kérésnek is explicit credential mód kell:

\begin{lstlisting}
const response = await fetch("https://api.example/api/account", {
  credentials: "include",
  headers: { Accept: "application/json" }
});
\end{lstlisting}

A \code{credentials: "include"} nem jogosultságot ad. Csak azt mondja a böngészőnek, hogy a megfelelő cookie-kat a CORS policy által megengedett requesthez elküldheti.

\section{CSRF JSON API-n}

Cookie-alapú sessionnél a böngésző a session cookie-t automatikusan kezeli, ezért állapotmódosító JSON requestnél ugyanúgy CSRF-védelem kell, mint HTML formnál.

Form POST esetén a token body mező:

\begin{lstlisting}
_csrf=<token>
\end{lstlisting}

JSON POST esetén header:

\begin{lstlisting}
X-CSRF-Token: <token>
\end{lstlisting}

Külön frontend originhez készíthető autentikált token endpoint:

\begin{lstlisting}[language=RWLang]
page fn csrfApi(ctx: PageContext) -> Result<Json, PageError> {
    return Ok(json(csrfToken));
}

route csrfApi GET "/api/csrf" auth user => csrfApi;
\end{lstlisting}

A token ugyanahhoz a sessionhöz tartozik. A CORS allowlist önmagában nem helyettesíti ezt az ellenőrzést.

\section{Teljes böngészős POST folyamat}

Egy külön \code{frontend.example} originről futó session-auth frontend tipikus lépései:

\begin{enumerate}
  \item a felhasználó autentikálódik, a böngésző megkapja a session cookie-t;
  \item a frontend lekéri a sessionhöz tartozó CSRF tokent;
  \item a POST kérés \code{credentials: "include"} módban megy;
  \item a request JSON bodyja megfelel a route zárt typed sémájának;
  \item az \code{X-CSRF-Token} header a session tokenjét tartalmazza;
  \item a runtime ellenőrzi az origint, authot, CSRF-et és input-validációt;
  \item a handler csak ezután hajt végre üzleti műveletet.
\end{enumerate}

Példa kliens:

\begin{lstlisting}
const csrf = await fetch("https://api.example/api/csrf", {
  credentials: "include",
  headers: { Accept: "application/json" }
}).then(async r => {
  if (!r.ok) throw new Error(`CSRF HTTP ${r.status}`);
  return r.json();
});

const result = await fetch("https://api.example/api/echo", {
  method: "POST",
  credentials: "include",
  headers: {
    "Content-Type": "application/json",
    Accept: "application/json",
    "X-CSRF-Token": csrf
  },
  body: JSON.stringify({ name: "Alice", age: 42, active: true })
}).then(async r => {
  if (!r.ok) throw new Error(`API HTTP ${r.status}`);
  return r.json();
});
\end{lstlisting}

\section{Reverse proxy és origin}

Apache vagy Nginx mögött különösen fontos, hogy a public scheme/host és a trusted proxy konfiguráció összhangban legyen. Az alkalmazás ne a kliens által szabadon küldött forwarded headerekből találja ki, hogy milyen origin tekinthető biztonságosnak. A trusted proxy boundary alapját a telepítési fejezet mutatja be; a teljes \code{server.toml} kulcsreferencia a könyv végén található.

A CORS allowlistbe mindig a \textbf{böngésző által látott frontend origint} írd, például:

\begin{lstlisting}
https://frontend.example
\end{lstlisting}

nem pedig a belső proxy vagy container címet.

\section{Mit ne tegyél?}

\begin{itemize}
  \item Ne engedélyezz CORS-t csak azért, mert ``API-ról van szó''.
  \item Ne tükrözd vissza tetszőleges request \code{Origin} értékét.
  \item Ne tekintsd a CORS-t auth- vagy CSRF-mechanizmusnak.
  \item Ne küldj session credentialt HTTP-n productionben.
  \item Ne keverd ugyanazon POST route-on a form, JSON és upload body módot.
  \item Ne hagyd korlátlanul a listázó query paramétereit; az API input is validált input.
  \item Ne építs klienslogikát HTML login redirect köré JSON route esetén; kezeld a 401/403 státuszt.
\end{itemize}

\section{Melyik architektúrát válaszd?}

\begin{center}
\begin{tabularx}{\textwidth}{@{}L{0.27\textwidth}Y@{}}
\toprule
\textbf{Helyzet} & \textbf{Ajánlott modell} \\
\midrule
szerver-renderelt oldal + kisebb JS & same-origin HTML + JSON API \\
azonos domainen futó SPA & same-origin API, CORS nélkül \\
külön publikus frontend, nincs session & exact-origin CORS credentials nélkül \\
külön frontend + session auth & exact-origin credentialed CORS + HTTPS + CSRF \\
server-to-server kliens & CORS nem releváns; az auth és hálózati policy a lényeg \\
\bottomrule
\end{tabularx}
\end{center}

\section{Kipróbálható példa}

A repository teljes mintája: \path{examples/json-api/app.rw}.

Elsőként ugyanazon originről érdemes kipróbálni. Ezután csak akkor vezess be külön frontend origint és CORS-t, ha a deployment valóban megköveteli. Így a fejlesztés során egyértelmű marad, mely problémát oldja meg a JSON API, és melyiket a böngésző cross-origin policy-ja.
